\chapter{Neural Networks and Their Limitations on Relational Data}
\label{ch:nn-limitations}
\minitoc

The preceding chapters established what graphs are and the tasks defined on
them. Before constructing a model that operates on graphs, it is necessary to
examine the standard tool of deep learning, the multilayer perceptron.
This chapter first assembles the multilayer
perceptron and the procedure by which it is trained, so that it stands as a
complete and well-defined object. It then attempts to apply this model to a
graph and shows that the attempt fails. The reasons for that failure are not
incidental: they identify precisely the properties a model must possess to
operate on graphs and motivate the developments 
presented in the following chapters.

\section{The Multilayer Perceptron}
\label{sec:mlp}

The \textit{multilayer perceptron} (MLP) is the elementary feedforward neural
network. It maps a fixed-size input vector to an
output vector through a sequence of layers, each applying a linear
transformation followed by a nonlinear function~\cite{bishop2024}.

A single layer takes a vector and produces another vector. Writing
$\mathbf{h}^{(l-1)}$ for the input to layer $l$ and $\mathbf{h}^{(l)}$ for its
output, the layer computes
\begin{equation}
  \mathbf{h}^{(l)} = \sigma\!\left(
    \mathbf{W}^{(l)} \mathbf{h}^{(l-1)} + \mathbf{b}^{(l)}
  \right),
  \label{eq:mlp-layer}
\end{equation}
where $\mathbf{W}^{(l)}$ is a learnable weight matrix, $\mathbf{b}^{(l)}$ a
learnable bias vector, and $\sigma$ the \textit{activation function}, 
a fixed nonlinear function applied elementwise. The index $l = 1, \ldots, L$
runs over the $L$ layers of the network. The weight matrix and 
bias vector of each layer are the \textit{parameters} of the model,
collectively denoted by $\theta$, and are the quantities adjusted during training.

The activation function $\sigma$ is what gives the network its expressive
power. Without it, the composition of linear layers would itself be a single
linear transformation, no matter how many layers were stacked. A nonlinear
$\sigma$ — for instance the rectified linear unit, $\sigma(x) = \max(0, x)$ —
allows the network to represent nonlinear relationships between input and
output~\cite{goodfellow2016}.

A multilayer perceptron is the composition of several such layers. Taking the
input vector $\mathbf{x}$ as the layer-zero representation,
$\mathbf{h}^{(0)} = \mathbf{x}$, the network applies the layers in sequence,
\begin{equation}
  f(\mathbf{x}) = f^{(L)} \circ \cdots \circ f^{(1)}(\mathbf{x}),
  \label{eq:mlp-composition}
\end{equation}
where each $f^{(l)}$ denotes the transformation of layer $l$ defined
in~\eqref{eq:mlp-layer}, and the output of the final layer is the prediction of
the network. The layers between the input and the output are the
\textit{hidden layers}, and the number of layers $L$ is the \textit{depth} of
the network. The map $f$ takes a fixed-size vector
$\mathbf{x} \in \mathbb{R}^d$ and returns an output of fixed size, determined
by the dimensions of the weight matrices (Figure~\ref{fig:mlp}).

\begin{figure}
  \centering
  \includegraphics{mlp}
  \caption[The multilayer perceptron.]{%
    A multilayer perceptron mapping a fixed-size input vector
    $\mathbf{x} \in \mathbb{R}^d$ through two hidden layers to an output. Each
    layer applies a learnable linear transformation followed by an elementwise
    nonlinearity. The dimension of the input is
    fixed in advance by the architecture: every input the network processes
    must be a vector of exactly $d$ components.
  }
  \label{fig:mlp}
\end{figure}

This last property is essential to what follows. The architecture fixes the
size of the input in advance: the product $\mathbf{W}^{(1)} \mathbf{x}$
is defined only if the input vector $\mathbf{x}$ has the required 
dimension, and the network assumes that its components are arranged 
in a definite order. A multilayer perceptron is a function on $\mathbb{R}^d$ 
for a fixed $d$, and on nothing else.

\section{Training by Gradient Descent}
\label{sec:training}
 
The parameters $\theta$ are not set by hand but learned from data, by adjusting
them so that the network's predictions match ground-truth targets. This requires a
measure of how wrong the predictions are and a procedure for reducing it.
 
The measure is a \textit{loss function}. Given a dataset of $n$ examples with
inputs $\mathbf{x}_i$ and targets $y_i$, the loss aggregates the discrepancy
between each prediction $f(\mathbf{x}_i)$ and its target into a single scalar,
\begin{equation}
  \mathcal{L}(\theta) = \frac{1}{n} \sum_{i=1}^{n}
    \ell\!\left(f(\mathbf{x}_i), y_i\right),
  \label{eq:loss}
\end{equation}
where $\ell$ is a per-example loss, such as the squared error for regression
or the cross-entropy for classification. The loss is a function of the
parameters $\theta$, since the prediction $f(\mathbf{x}_i)$ depends on them;
training is the search for the $\theta$ that minimizes $\mathcal{L}(\theta)$.
 
That search is carried out by \textit{gradient descent}. The gradient
$\nabla_\theta \mathcal{L}$ points in the direction of steepest increase of the
loss; moving the parameters in the opposite direction decreases it. Each step
updates the parameters by
\begin{equation}
  \theta \leftarrow \theta - \eta \, \nabla_\theta \mathcal{L}(\theta),
  \label{eq:gradient-descent}
\end{equation}
where the scalar $\eta > 0$ is the \textit{learning rate}, controlling the size
of each step. The gradient itself is computed by \textit{backpropagation}, the
application of the chain rule across the layers of the
network~\cite{goodfellow2016}; it yields the derivative of the loss with
respect to every parameter, and is the algorithm that makes training deep
networks practical.
 
Computing the gradient in~\eqref{eq:gradient-descent} over the full dataset is
costly when $n$ is large, since it requires a pass over every example for a
single update. \textit{Stochastic gradient descent} replaces the full-dataset
gradient with an estimate computed on a small random subset of the data, a
\textit{minibatch} $\mathcal{B}$. The method was originally defined for a single
example per step — the source of the name, as each step follows a noisy,
\textit{stochastic} estimate of the gradient rather than the exact one — but is
generally applied over a minibatch, which averages the per-example gradients,
\begin{equation}
  \nabla_\theta \mathcal{L}(\theta) \approx
    \frac{1}{\lvert \mathcal{B} \rvert} \sum_{i \in \mathcal{B}}
    \nabla_\theta \, \ell\!\left(f(\mathbf{x}_i), y_i\right).
  \label{eq:sgd}
\end{equation}
Under uniform random sampling of examples, the minibatch gradient is an 
unbiased estimate of the full gradient, and each update is far cheaper, 
so many more updates are performed for the same computational cost. 
Stochastic gradient descent is the basis of essentially all
modern neural network training.
 
The optimizer used throughout the implementations of this work is
\textit{Adam}~\cite{kingma2015}, a method in the stochastic gradient descent
family that adapts the effective step size for each parameter individually,
using running estimates of the gradient. Its internal mechanics are not required for the
developments that follow, and the reader is referred to the original work for
them. What matters is that the multilayer perceptron, with its parameters, loss, and training
procedure, now stands as a complete and trainable object. The remainder of the
chapter asks whether that object can be applied to a graph.

\section{Why the Multilayer Perceptron Fails on Graphs}
\label{sec:mlp-fails}

The multilayer perceptron assembled above is a function on a fixed-size vector.
A graph is not a vector. To apply the perceptron to a graph, the graph must
first be converted into a vector of fixed size, and the difficulty is that no
such conversion preserves what the graph encodes. This section attempts the
conversion directly and examines what goes wrong; the failures that surface are
the substance of the chapter.

Consider a node-level task: predicting a property of each node of a graph, as
in the node classification setting of Chapter~\ref{ch:applications}. There are
two natural ways to present such a problem to a multilayer perceptron, and it
is instructive that both fail, for different reasons.

The first approach attempts to encode the whole graph as a single input vector
and feed it to the network. The structure of a graph on $n$ nodes is carried by
its adjacency matrix $\mathbf{A} \in \mathbb{R}^{n \times n}$, and its node
attributes by the feature matrix $\mathbf{X} \in \mathbb{R}^{n \times d}$. To
form a single input vector, these matrices are flattened — their entries read
out in order and laid end to end — and concatenated. This at once runs into two
obstructions.

The first obstruction concerns size. The flattened adjacency matrix has $n^2$ entries and
the flattened feature matrix has $n \times d$; the length of the resulting
input vector therefore grows with the number of nodes and differs from one
graph to another. But a multilayer perceptron, by~\eqref{eq:mlp-layer}, admits
inputs of one fixed dimension only. A network built for a graph of one size
cannot accept a graph of another. The same obstruction appears within a single
graph. A finer attempt would process each node together with its neighbors,
rather than flattening the whole graph at once; but nodes differ in how many
neighbors they have, so this neighborhood-based input again has no fixed size
(Figure~\ref{fig:variable-neighborhood}). This is the problem of \textit{variable
size}: neither across graphs nor across the nodes of one graph is the number of
entities fixed, whereas the multilayer perceptron demands a
fixed-dimensional input.

\begin{figure}
  \centering
  \includegraphics{variable-neighborhood}
  \caption[Neighborhoods of varying size.]{%
    Two nodes of the same graph with neighborhoods of different sizes: the
    node on the left has three neighbors, the node on the right has six. A model
    that processes a node together with its neighbors must accept inputs of
    varying size, which a multilayer perceptron cannot.
  }
  \label{fig:variable-neighborhood}
\end{figure}

The second obstruction concerns order. As established in
Chapter~\ref{ch:graph-theory}, the nodes of a graph carry no intrinsic
ordering; an indexing must be chosen to write down the adjacency matrix, and
any of the $n!$ possible indexings describes the same graph. The two adjacency
matrices of a graph under two different node orderings are different matrices,
and flattening them yields two different vectors. A multilayer perceptron
would receive these as two distinct inputs and could process them differently,
even though they represent the same graph
(Figure~\ref{fig:relabeling}). The model would thus depend on an arbitrary
choice that carries no information — the labeling of the nodes — and would have
no means of recognizing that the choice is arbitrary. This is the problem of
\textit{node ordering}: the flattened representation imposes an order that the
graph does not have, and the multilayer perceptron has no way to disregard it.

\begin{figure*}[t]
  \centering
  \includegraphics[width=\textwidth]{relabeling}
  \caption[The same graph under two node orderings.]{%
    A single graph drawn twice with its nodes indexed in two different orders.
    Each ordering produces a different adjacency matrix, and flattening each
    matrix produces a different input vector. A multilayer perceptron receives
    the two vectors as unrelated inputs, although they encode the same graph.
    Neither ordering is privileged: the graph is the same object in both.
  }
  \label{fig:relabeling}
\end{figure*}

Faced with these obstructions, one might abandon the adjacency matrix
altogether and take the second approach: present each node to the network on
its own, using only its feature vector $\mathbf{x}_i$, and predict the node's
property from that. This input has a fixed size $d$, and the ordering of the
nodes no longer enters, since each node is treated separately. Both previous
obstructions vanish — but at a cost that defeats the purpose. With the
adjacency matrix discarded, the structure of the graph never enters the
computation at all. The network processes each node as an isolated example. Yet
the tasks of Chapter~\ref{ch:applications} are relational by nature: in a
citation network, a paper's topic is informed by the topics of the papers it
cites, and a model that sees a node in isolation discards exactly this
evidence. This is the problem of \textit{missing relational structure}: the
only input that fits a multilayer perceptron is one from which the relations
have been removed.

The two approaches exhaust the natural options, and they fail in opposite
ways. Including the adjacency matrix preserves the structure but breaks on
variable size and node ordering; excluding it restores a fixed, order-free
input but discards the structure entirely. There is no middle ground: the
multilayer perceptron cannot be given a graph without either confronting a
representation it cannot accept or being handed one from which the graph has
been erased.

These failures reveal that the multilayer perceptron carries the wrong
assumptions for relational data. Its \textit{inductive bias} is that examples
are independent and of fixed size, an assumption appropriate to vectors but
false for graphs, whose entities are connected and whose representation admits
no canonical order. What the analysis exposes is the assumptions a
graph model must instead possess: it must accept inputs of varying size, it
must be \textit{invariant} to the ordering of the nodes, so that relabeling
leaves its predictions unchanged, and it must let the structure of the graph
participate in the computation. These three requirements, arrived at here by
watching the multilayer perceptron fail, are the foundation on which the following
chapters build a neural network suited to graphs.